\clearpage
\begingroup
\tcbset{myboxstyle/.append style={breakable=false}}
\everymath{\textstyle}
\subsection{Случайный вектор и совместная функция распределения}

\begin{definition}[Случайный вектор]
    \textbf{Случайным вектором} называется вектор
    $\vec\xi=(\xi_1,\ldots,\xi_n)$, компоненты которого являются
    случайными величинами на одном вероятностном пространстве
    $(\Omega,\mathcal F,\Pr)$.
\end{definition}

\begin{definition}[Совместная функция распределения]
    \textbf{Функцией распределения случайного вектора} $\vec\xi$
    называется функция
    \[
        F_{\vec\xi}(x_1,\ldots,x_n)
        =\Pr\{\xi_1\leq x_1,\ldots,\xi_n\leq x_n\}.
    \]
    Иными словами,
    \[
        F_{\vec\xi}(\vec x)
        =\Pr\left(\bigcap_{i=1}^{n}\{\omega:\xi_i(\omega)\leq x_i\}\right).
    \]
    Она называется также \textbf{совместной функцией распределения}
    величин $\xi_1,\ldots,\xi_n$.
\end{definition}

\begin{theorem}[Свойства совместной функции распределения]
    \begin{enumerate}
        \item $0\leq F_{\vec\xi}(\vec x)\leq1$.
        \item Предельный переход по одной координате даёт функцию
        распределения остальных координат. Для пары $(\xi,\eta)$:
        \[
            \begin{aligned}
                \lim_{x\to+\infty}F_{\xi,\eta}(x,y)&=F_\eta(y),&
                \lim_{x\to-\infty}F_{\xi,\eta}(x,y)&=0,\\
                \lim_{y\to+\infty}F_{\xi,\eta}(x,y)&=F_\xi(x),&
                \lim_{y\to-\infty}F_{\xi,\eta}(x,y)&=0.
            \end{aligned}
        \]
        \item $F_{\vec\xi}$ не убывает по каждой координате.
        \item Если $\vec x^{\,(k)}\downarrow\vec x^{\,(0)}$
        покоординатно, то
        \[
            F_{\vec\xi}\bigl(\vec x^{\,(k)}\bigr)
            \downarrow F_{\vec\xi}\bigl(\vec x^{\,(0)}\bigr).
        \]
    \end{enumerate}
\end{theorem}

\newpage
\subsection{Независимость случайных величин}

События $A$ и $B$ независимы, если $\Pr(A\cap B)=\Pr(A)\Pr(B)$.

\begin{definition}
    Случайные величины $\xi$ и $\eta$ называются \textbf{независимыми}, если
    \[
        F_{\xi,\eta}(x,y)=F_\xi(x)F_\eta(y)
        \qquad\forall x,y\in\R.
    \]
\end{definition}

\begin{theorem}[Эквивалентное определение через события]
    Величины $\xi$ и $\eta$ независимы тогда и только тогда, когда
    для любых $A,B\in\mathcal B(\R)$ события $\{\xi\in A\}$ и $\{\eta\in B\}$
    независимы, то есть
    \[
        \Pr\{\xi\in A,\eta\in B\}
        =\Pr\{\xi\in A\}\Pr\{\eta\in B\}.
    \]
\end{theorem}

\begin{Proof}[От независимости событий к разложению функции распределения]
    Подставим $A=(-\infty;x]$, $B=(-\infty;y]$. Тогда
    \[
        \begin{aligned}
            F_{\xi,\eta}(x,y)
            &=\Pr\{\xi\leq x,\eta\leq y\}\\
            &=\Pr\{\xi\leq x\}\Pr\{\eta\leq y\}
             =F_\xi(x)F_\eta(y).
        \end{aligned}
    \]
\end{Proof}

\begin{example}[События случайной величины]
    Пусть $\Omega=[0;1]$, $\mathcal F=\mathcal B([0;1])$, а $\Pr$ ---
    мера Лебега на $[0;1]$. Для $\xi(\omega)=\lfloor2\omega\rfloor$ имеем
    \[
        \xi(\omega)=
        \begin{cases}
            0,&0\leq\omega<\frac12,\\
            1,&\frac12\leq\omega<1,\\
            2,&\omega=1.
        \end{cases}
    \]
    Поэтому каждое событие $\{\xi\in B\}$ является объединением некоторых
    из трёх множеств
    \[
        [0;\tfrac12),\qquad[\tfrac12;1),\qquad\{1\}.
    \]
    Всего возможны восемь таких объединений, включая $\varnothing$ и $\Omega$.
\end{example}

\newpage
\subsection{Совместное дискретное распределение}

\begin{definition}
    Пусть $\xi$ и $\eta$ принимают значения в конечных или счётных
    множествах $\{x_i\}$ и $\{y_j\}$ с вероятностью $1$.
    Их \textbf{совместное распределение} задаётся числами
    \[
        p_{ij}=\Pr\{\xi=x_i,\eta=y_j\}.
    \]
    Вероятности отдельных значений равны
    \[
        p_i=\Pr\{\xi=x_i\}=\sum_jp_{ij},\qquad
        q_j=\Pr\{\eta=y_j\}=\sum_ip_{ij}.
    \]
    При этом $p_{ij}\geq0$ и $\sum_{i,j}p_{ij}=1$.
\end{definition}

\begin{theorem}[Критерий независимости]
    Дискретные случайные величины $\xi$ и $\eta$ независимы
    тогда и только тогда, когда
    \[
        p_{ij}=p_iq_j\qquad\forall i,j.
    \]
\end{theorem}

\begin{theorem}[Матричный критерий]
    Если множества значений обеих величин конечны, то для матрицы
    совместных вероятностей $P=(p_{ij})$
    \[
        \xi\text{ и }\eta\text{ независимы}
        \quad\Longleftrightarrow\quad\rk P=1.
    \]
    В случае независимости
    \[
        P=
        \begin{pmatrix}
            p_1q_1&\cdots&p_1q_s\\
            \vdots&\ddots&\vdots\\
            p_rq_1&\cdots&p_rq_s
        \end{pmatrix}
        =\begin{pmatrix}p_1\\\vdots\\p_r\end{pmatrix}
         \begin{pmatrix}q_1&\cdots&q_s\end{pmatrix}.
    \]
\end{theorem}

\newpage
\subsection{Совместная плотность распределения}

\begin{definition}
    Распределение пары $(\xi,\eta)$ называется \textbf{абсолютно непрерывным},
    если существует неотрицательная измеримая функция $p_{\xi,\eta}$,
    для которой
    \[
        F_{\xi,\eta}(x,y)
        =\int_{-\infty}^{y}\int_{-\infty}^{x}
            p_{\xi,\eta}(u,v)\,\diff u\,\diff v.
    \]
    Функция $p_{\xi,\eta}$ называется \textbf{совместной плотностью}.
    Она удовлетворяет условию
    \[
        \iint_{\R^2}p_{\xi,\eta}(u,v)\,\diff u\,\diff v=1.
    \]
\end{definition}

\begin{theorem}[Критерий независимости через плотности]
    Если пара $(\xi,\eta)$ имеет совместную плотность, то
    \[
        \xi\text{ и }\eta\text{ независимы}
        \quad\Longleftrightarrow\quad
        p_{\xi,\eta}(x,y)=p_\xi(x)p_\eta(y)
        \ \text{почти всюду}.
    \]
    Здесь почти всюду понимается относительно меры Лебега на $\R^2$.
\end{theorem}

\begin{Proof}[Необходимость]
    При независимости $F_{\xi,\eta}(x,y)=F_\xi(x)F_\eta(y)$.
    Дифференцируя по обеим переменным, получаем почти всюду
    \[
        p_{\xi,\eta}(x,y)
        =\frac{\partial^2F_{\xi,\eta}}{\partial x\,\partial y}(x,y)
        =F_\xi'(x)F_\eta'(y)=p_\xi(x)p_\eta(y).
    \]
\end{Proof}

\begin{Proof}[Достаточность]
    Если $p_{\xi,\eta}(u,v)=p_\xi(u)p_\eta(v)$ почти всюду, то
    \[
        \begin{aligned}
            F_{\xi,\eta}(x,y)
            &=\int_{-\infty}^{y}\int_{-\infty}^{x}
                p_\xi(u)p_\eta(v)\,\diff u\,\diff v\\
            &=\left(\int_{-\infty}^{x}p_\xi(u)\,\diff u\right)
              \left(\int_{-\infty}^{y}p_\eta(v)\,\diff v\right)\\
            &=F_\xi(x)F_\eta(y).
        \end{aligned}
    \]
\end{Proof}
\endgroup
